%% =====================================================================
%%  Unit 3 — What's your number?  (Les lettres et les nombres)
%%  Pistes : 3.01–3.04 dialogues · 3.05–3.08 vocabulaire · 3.09 grammaire
%%           3.10 prononciation · 3.11–3.13 exercices (textes en annexe)
%% =====================================================================

\newcommand{\otIllusUnIII}{\otIllusDuo{A B C/104/86/15/0.10,1 2 3/190/92/20/0.09,@/122/14/16/0.08,
  zero/150/8/11/0.08,twenty/196/22/12/0.07,X Y Z/100/50/12/0.07}{emma}{What's your number?}{lea}{917-555-0148}}

\unite[What's your\\number?]{3}{What's your number?}{Les lettres et les nombres}{otIllusUnIII}

\situation{Épisode 3}{Mercredi soir. Léa s'inscrit dans une salle de sport de Brooklyn : il faut
épeler son nom, donner son numéro de téléphone, son adresse e-mail… et son âge. Le lendemain, au
studio, Emma et Léa échangent leurs numéros. Et qui est ce jeune homme qui arrive avec Jack ?}

\begin{objectifs}
\item épeler votre nom et comprendre un mot épelé ;
\item compter de 0 à 100 ;
\item donner votre âge, votre numéro de téléphone et votre adresse e-mail ;
\item poser des questions avec \en{what}, \en{where}, \en{who}, \en{how} et \en{how old} ;
\item choisir entre \en{a} et \en{an}.
\end{objectifs}

\auprogramme
  {Épeler, compter, donner son âge, son numéro et son adresse e-mail}
  {Les mots interrogatifs\par\vspace{0.6mm}\en{a} ou \en{an} ?}
  {L'alphabet, les nombres de 0 à 100 et 16 expressions}
  {Les lettres qui piègent : E, I, G, J\ldots}
  {Numéros de téléphone et adresses e-mail}

\semaine{casque/otNavy/Dialogues/Écouter et lire les deux scènes,
  bulle/otGreenDark/Vocabulaire/L'alphabet et les nombres,
  blocs/otViolet/Grammaire/Mots interrogatifs et a ou an,
  sons/otOrangeDark/Sons et culture/Les lettres qui piègent,
  crayon/otNavy/Exercices/12 exercices corrigés,
  coche/otGreenDark/Bilan/L'essentiel et le test express}

\clearpage
%% ---------------------------------------------------------------------
\rubrique{ecouter}{À la salle de sport}

\begin{dialogue}{1}{How do you spell that?}{Mercredi, 19 heures. Léa s'inscrit à \en{Fit Brooklyn},
  une salle de sport près de chez elle. L'employée remplit sa fiche.}{3.01}{3.02}
\dit{employee}{Hi! Welcome to Fit Brooklyn. What's your name, please?}{Bonjour ! Bienvenue chez Fit Brooklyn. Votre nom, s'il vous plaît ?}
\dit{lea}{Léa Martin.}{Léa Martin.}
\dit{employee}{\hl{How do you spell} your \hl{last name}?\note{1}}{Comment s'écrit votre nom de famille ?}
\dit{lea}{Martin: M, A, R, T, I, N.}{Martin : M, A, R, T, I, N.}
\dit{employee}{And your \hl{first name}?}{Et votre prénom ?}
\dit{lea}{Léa: L, E, A.}{Léa : L, E, A.}
\dit{employee}{Thank you. \hl{What's your phone number}?}{Merci. Quel est votre numéro de téléphone ?}
\dit[It's nine one seven, five five five, oh one four eight.]{lea}{It's 917-555-0148.\note{2}}{C'est le 917 555 0148.}
\dit{employee}{And your \hl{email address}?}{Et votre adresse e-mail ?}
\dit[It's Léa dot Martin, at example dot com.]{lea}{It's lea.martin@example.com.\note{3}}{C'est lea.martin@example.com.}
\dit{employee}{OK. And \hl{how old are you}?}{D'accord. Et quel âge avez-vous ?}
\dit{lea}{\hl{I'm twenty-seven}.\note{4}}{J'ai vingt-sept ans.}
\dit{employee}{Perfect. Here's your card. Welcome to the club!}{Parfait. Voici votre carte. Bienvenue au club !}
\end{dialogue}

\begin{comprendre}
\pointnote{1} \en{\textbf{How do you spell…?}} « Comment épelez-vous… ? ». \en{to spell} = épeler.
  \en{last name} = nom de famille ; \en{first name} = prénom.
\pointnote{2} \en{\textbf{917-555-0148}} On lit un numéro de téléphone \textbf{chiffre par chiffre},
  en marquant une pause entre les groupes : \en{nine one seven, five five five, oh one four eight}.
  Le zéro se dit le plus souvent \en{oh}, comme la lettre O.
\pointnote{3} \en{\textbf{lea.martin@example.com}} Dans une adresse e-mail, \en{@} se dit
  \en{at} et le point se dit \en{dot}.
\pointnote{4} \en{\textbf{I'm twenty-seven.}} En anglais, on \textbf{est} un âge : « je suis
  vingt-sept ». On peut ajouter \en{years old} : \en{I'm twenty-seven years old.} Mais jamais
  \en{I have twenty-seven years} !
\end{comprendre}

\begin{dialogue}{2}{What's your number?}{Jeudi matin, au studio. Emma veut le numéro de Léa.}{3.03}{3.04}
\dit{emma}{Léa, what's your number?}{Léa, c'est quoi ton numéro ?}
\dit[It's nine one seven, five five five, oh one four eight.]{lea}{It's 917-555-0148.}{C'est le 917 555 0148.}
\dit{emma}{\hl{Sorry? One more time}, please.\note{5}}{Pardon ? Encore une fois, s'il te plaît.}
\dit{lea}{Nine one seven, five five five, oh one four eight.}{Neuf un sept, cinq cinq cinq, zéro un quatre huit.}
\dit[Got it! And my number is six four six, five five five, oh one nine three.]{emma}{Got it! And my number is 646-555-0193.}{C'est noté ! Et mon numéro, c'est le 646 555 0193.}
\dit{lea}{Thanks. \hl{By the way}, how old are you, Emma?\note{6}}{Merci. Au fait, quel âge as-tu, Emma ?}
\dit{emma}{I'm twenty-nine. And Jack is thirty.}{J'ai vingt-neuf ans. Et Jack a trente ans.}
\dit{lea}{Thirty? Really? \hl{Who's that}, with Jack?\note{7}}{Trente ans ? Vraiment ? Qui est-ce, avec Jack ?}
\dit{emma}{\hl{That's} Tom. He's new here. He's nineteen!}{C'est Tom. Il est nouveau ici. Il a dix-neuf ans !}
\end{dialogue}

\begin{comprendre}
\pointnote{5} \en{\textbf{Sorry? One more time, please.}} \en{Sorry?}, avec la voix qui monte,
  veut dire « Pardon ? » quand on n'a pas compris. \en{One more time} = « encore une fois ».
  \en{Got it!} = « C'est noté ! », « Compris ! ».
\pointnote{6} \en{\textbf{By the way}} « Au fait », pour changer de sujet.
\pointnote{7} \en{\textbf{Who's that?}} « Qui est-ce ? ». \en{Who's} = \en{Who is} ;
  \en{That's} = \en{That is}. \en{new} = nouveau, nouvelle.
\end{comprendre}

%% ---------------------------------------------------------------------
\rubrique{mots}{Les lettres, les nombres et les mots de l'unité}

Au programme : l'alphabet, les nombres de 0 à 100 et 16 expressions pour échanger ses
coordonnées. Attention aux lettres dont le nom ressemble à une autre lettre en français :
le \en{E} anglais se dit comme notre « i », et le \en{I} se dit « aï ».

\begin{grille}{3.05}{L'alphabet}{7}
\grcase{A}{eɪ}{} \grcase{B}{biː}{} \grcase{C}{siː}{} \grcase{D}{diː}{} \grcase{E}{iː}{} \grcase{F}{ef}{} \grcase{G}{dʒiː}{}
\grcase{H}{eɪtʃ}{} \grcase{I}{aɪ}{} \grcase{J}{dʒeɪ}{} \grcase{K}{keɪ}{} \grcase{L}{el}{} \grcase{M}{em}{} \grcase{N}{en}{}
\grcase{O}{oʊ}{} \grcase{P}{piː}{} \grcase{Q}{kjuː}{} \grcase{R}{ɑːr}{} \grcase{S}{es}{} \grcase{T}{tiː}{} \grcase{U}{juː}{}
\grcase{V}{viː}{} \grcase{W}{ˈdʌbəljuː}{} \grcase{X}{eks}{} \grcase{Y}{waɪ}{} \grcase{Z}{ziː}{}
\end{grille}

\begin{grille}{3.06}{Les nombres de 0 à 20}{7}
\grcase{0}{ˈzɪroʊ}{zero} \grcase{1}{wʌn}{one} \grcase{2}{tuː}{two} \grcase{3}{θriː}{three}
\grcase{4}{fɔːr}{four} \grcase{5}{faɪv}{five} \grcase{6}{sɪks}{six}
\grcase{7}{ˈsevən}{seven} \grcase{8}{eɪt}{eight} \grcase{9}{naɪn}{nine} \grcase{10}{ten}{ten}
\grcase{11}{ɪˈlevən}{eleven} \grcase{12}{twelv}{twelve} \grcase{13}{θɝːˈtiːn}{thirteen}
\grcase{14}{ˌfɔːrˈtiːn}{fourteen} \grcase{15}{ˌfɪfˈtiːn}{fifteen} \grcase{16}{ˌsɪkˈstiːn}{sixteen}
\grcase{17}{ˌsevənˈtiːn}{seventeen} \grcase{18}{ˌeɪˈtiːn}{eighteen} \grcase{19}{ˌnaɪnˈtiːn}{nineteen}
\grcase{20}{ˈtwenti}{twenty}
\end{grille}

\begin{grille}{3.07}{Les dizaines, jusqu'à 100}{6}
\grcase{30}{ˈθɝːti}{thirty} \grcase{40}{ˈfɔːrti}{forty} \grcase{50}{ˈfɪfti}{fifty}
\grcase{60}{ˈsɪksti}{sixty} \grcase{70}{ˈsevənti}{seventy} \grcase{80}{ˈeɪti}{eighty}
\grcase{90}{ˈnaɪnti}{ninety} \grcase{100}{ə ˈhʌndrəd}{a hundred} \grcase{21}{ˌtwentiˈwʌn}{twenty-one}
\grcase{35}{ˌθɝːtiˈfaɪv}{thirty-five} \grcase{48}{ˌfɔːrtiˈeɪt}{forty-eight} \grcase{99}{ˌnaɪntiˈnaɪn}{ninety-nine}
\end{grille}

\begin{astuce}[Former les nombres de 21 à 99]
Dizaine + trait d'union + unité : \en{twenty-one}, \en{thirty-five}, \en{ninety-nine}. Pas de
« et » comme dans « vingt et un », et aucune exception du type « soixante-dix » : 70 se dit
simplement \en{seventy}, 80 \en{eighty} et 90 \en{ninety}.
\end{astuce}

\begin{vocabulaire}{Échanger ses coordonnées}{3.08}
\categorie{Questions utiles}
\mot{How do you spell it?}{haʊ duː juː ˈspel ɪt}{Comment ça s'écrit ?}
\mot{What's your phone number?}{ˌwʌts jʊr ˈfoʊn ˌnʌmbɚ}{Quel est ton numéro de téléphone ?}
\mot{What's your email address?}{ˌwʌts jʊr ˈiːmeɪl ˌædres}{Quelle est ton adresse e-mail ?}
\mot{How old are you?}{ˌhaʊ ˈoʊld ɑːr juː}{Quel âge as-tu ? / Quel âge avez-vous ?}
\mot{I'm twenty-seven.}{aɪm ˌtwentiˈsevən}{J'ai vingt-sept ans.}
\mot{Who's that?}{ˌhuːz ˈðæt}{Qui est-ce ?}
\mot{That's Jack.}{ˌðæts ˈdʒæk}{C'est Jack.}
\categorie{Les mots des coordonnées}
\mot{first name}{ˈfɝːst ˌneɪm}{prénom}
\mot{last name}{ˈlæst ˌneɪm}{nom de famille}
\mot{a phone number}{ə ˈfoʊn ˌnʌmbɚ}{un numéro de téléphone}
\mot{an email address}{ən ˈiːmeɪl ˌædres}{une adresse e-mail}
\mot{at}{æt}{arobase (@)}
\mot{dot}{dɑːt}{point (dans une adresse)}
\categorie{Pour se faire répéter}
\mot{Sorry?}{ˈsɑːri}{Pardon ?}
\mot{One more time, please.}{ˌwʌn mɔːr ˈtaɪm pliːz}{Encore une fois, s'il vous plaît.}
\mot{Got it!}{ˈɡɑːt ɪt}{C'est noté ! Compris !}
\end{vocabulaire}

%% ---------------------------------------------------------------------
\rubrique{grammaire}{Les mots interrogatifs, a ou an}

\otIntertitre[otViolet]{1 · Les mots interrogatifs}
Le mot interrogatif se place \textbf{en tête} de la question ; viennent ensuite le verbe \en{be}
et le sujet. À l'oral, \en{is} se contracte : \en{What's}, \en{Where's}, \en{Who's}.

\begin{center}\otTableau
\begin{tabular}{@{}>{\enOn\bfseries}l<{\enOff}@{\hspace{6mm}}l@{\hspace{6mm}}>{\enOn}l<{\enOff}@{\hspace{6mm}}>{\enOn}l<{\enOff}@{}}
\otTH{Mot} & \otTH{Sens} & \otTH{Question} & \otTH{Réponse}\\\midrule
What     & quel, quelle, quoi & What's your name?       & My name is Léa.\\
Where    & où                 & Where are you from?     & I'm from France.\\
Who      & qui                & Who's that?             & That's Jack.\\
How      & comment            & How are you?            & I'm fine, thanks.\\
How old  & quel âge           & How old are you?        & I'm twenty-seven.\\\bottomrule
\end{tabular}\end{center}

\blocaudio{3.09}{Écoutez et répétez}{Six questions et leur réponse. Répétez la question, puis la
réponse, pendant les pauses.}

\otIntertitre[otViolet]{2 · \en{a} ou \en{an} ?}
\en{a} et \en{an} veulent dire « un » ou « une ». On choisit selon le \textbf{son} qui suit, et
non selon la lettre :

\begin{paire}
\begin{regle}[\en{a} devant un son consonne]
\exen{a phone number}{}
\exen{a city}{}
\exen{a beginner}{}
\exen{a university}{\en{u} se dit \phon{juː} : son consonne}
\end{regle}
\begin{regle}[\en{an} devant un son voyelle]
\exen{an email address}{}
\exen{an old city}{}
\exen{an American}{}
\exen{an hour}{le \en{h} est muet : son voyelle}
\end{regle}
\end{paire}

\begin{attention}[Les pièges des francophones]
\fauxjuste{I have twenty-seven years.}{I'm twenty-seven.}
\fauxjuste{How old have you?}{How old are you?}
\fauxjuste{a email address}{an email address}
\fauxjuste{an university}{a university}
\fauxjuste{fourty}{forty}
\par\vspace{1mm}{\fontsize{10}{12.6}\selectfont \en{four} et \en{fourteen} gardent le \en{u}, mais
pas \en{forty}. Et l'on écrit \en{ninety}, avec un \en{e}.\par}
\end{attention}

%% ---------------------------------------------------------------------
\rubrique[Les lettres qui piègent]{sons}{Les lettres qui piègent}

Le nom des lettres anglaises ressemble souvent à celui d'une \textbf{autre} lettre française.
C'est la principale source d'erreurs quand on épelle au téléphone. Retenez surtout ces deux
couples, inversés par rapport au français : \en{E} \phon{iː} et \en{I} \phon{aɪ} ;
\en{G} \phon{dʒiː} et \en{J} \phon{dʒeɪ}.

\begin{paires}{3.10}{Écoutez les deux lettres de chaque paire, puis répétez-les en exagérant la
différence.}
\sonpaire{E}{iː}{comme le « i » de « vie »}{I}{aɪ}{comme « aïe »}
\sonpaire{G}{dʒiː}{« dji »}{J}{dʒeɪ}{« djéï »}
\sonpaire{A}{eɪ}{« éï »}{R}{ɑːr}{« ar », avec le r anglais}
\sonpaire{I}{aɪ}{« aï »}{Y}{waɪ}{« ouaï »}
\sonpaire{K}{keɪ}{« kéï »}{Q}{kjuː}{« kiou »}
\sonpaire{B}{biː}{« bi »}{V}{viː}{« vi »}
\sonpaire{U}{juː}{« iou »}{W}{ˈdʌbəljuː}{« double iou »}
\sonpaire{S}{es}{« ès »}{X}{eks}{« èks »}
\end{paires}

\begin{astuce}[Épeler sans erreur]
Pour une lettre doublée, les Américains répètent la lettre (\en{B, R, O, O, K, S}) ; les
Britanniques disent souvent \en{double} : \en{double O}. Au téléphone, on précise une lettre avec
un mot : \en{M as in Mike}, \en{N as in November} (« M comme… »).
\end{astuce}

%% ---------------------------------------------------------------------
\rubrique{culture}{Numéros et adresses}

\begin{paire}
\begin{culture}[Les numéros de téléphone]
Aux États-Unis, un numéro compte dix chiffres : un indicatif de zone (\en{area code}) de trois
chiffres, puis sept chiffres, notés \en{(917) 555-0148}. 212 et 646 sont des indicatifs de
Manhattan, 718 et 347 de Brooklyn et des autres quartiers.
\end{culture}
\begin{culture}[Le 555 du cinéma]
Dans les films et les séries, les numéros commencent souvent par 555 : ceux de 555-0100 à
555-0199 sont réservés à la fiction et n'appartiennent à personne. Tous les numéros de ce livre
sont choisis parmi eux.
\end{culture}
\end{paire}

\begin{paire}
\begin{culture}[Les urgences]
Aux États-Unis et au Canada, on compose le \textbf{911} (\en{nine one one}) ; au Royaume-Uni, le
\textbf{999}. Le 112 européen fonctionne aussi au Royaume-Uni.
\end{culture}
\begin{culture}[Lire une adresse e-mail]
\en{@} : \en{at} · \en{.} : \en{dot} · \en{-} : \en{dash} (ou \en{hyphen}) · \en{\_} :
\en{underscore}. Et la dernière lettre de l'alphabet se dit \en{zee} aux États-Unis, mais
\en{zed} au Royaume-Uni !
\end{culture}
\end{paire}

%% ---------------------------------------------------------------------
\rubrique{pratique}{Exercices}

\exercice[ecrire,niveau=1]{En toutes lettres}{Écrivez chaque nombre en lettres.}
\begin{questions}[2]
\item 7 → \trou[34mm]{seven}
\item 12 → \trou[34mm]{twelve}
\item 15 → \trou[34mm]{fifteen}
\item 20 → \trou[34mm]{twenty}
\item 33 → \trou[34mm]{thirty-three}
\item 48 → \trou[34mm]{forty-eight}
\item 50 → \trou[34mm]{fifty}
\item 99 → \trou[34mm]{ninety-nine}
\end{questions}

\exercice[lire,niveau=1]{En chiffres}{Écrivez chaque nombre en chiffres.}
\begin{questions}[3]
\item thirteen → \trou[12mm]{13}
\item thirty → \trou[12mm]{30}
\item eighty-one → \trou[12mm]{81}
\item sixty-six → \trou[12mm]{66}
\item fourteen → \trou[12mm]{14}
\item forty → \trou[12mm]{40}
\end{questions}

\exercice[ecrire,niveau=1]{\en{a} ou \en{an} ?}{Complétez avec \en{a} ou \en{an}.}
\begin{questions}[2]
\item \trou[12mm]{a} phone number
\item \trou[12mm]{an} email address
\item \trou[12mm]{a} big city
\item \trou[12mm]{an} hour
\item \trou[12mm]{a} beginner
\item \trou[12mm]{an} Italian
\item \trou[12mm]{a} university
\item \trou[12mm]{an} old city
\end{questions}

\exercice[ecrire,niveau=1]{Le bon mot interrogatif}{Complétez avec \en{What}, \en{Where},
\en{Who}, \en{How} ou \en{How old}.}
\begin{questions}
\item \trou[22mm]{What}'s your name? — It's Léa.
\item \trou[22mm]{Where} are you from? — I'm from France.
\item \trou[22mm]{How old} are you? — I'm twenty-seven.
\item \trou[22mm]{How} are you? — I'm fine, thanks.
\item \trou[22mm]{Who}'s that? — That's Mr.\ Carter.
\item \trou[22mm]{What}'s your phone number? — It's 917-555-0148.
\end{questions}

\exercice[ecrire,niveau=2]{Dans le bon ordre}{Remettez les mots dans l'ordre.}
\begin{questions}
\item \etiq{old} \etiq{are} \etiq{how} \etiq{you} \etiq{?}\hfill\trou[66mm]{How old are you?}
\item \etiq{spell} \etiq{you} \etiq{do} \etiq{how} \etiq{it} \etiq{?}\hfill\trou[66mm]{How do you spell it?}
\item \etiq{number} \etiq{phone} \etiq{what's} \etiq{your} \etiq{?}\hfill\trou[66mm]{What's your phone number?}
\item \etiq{twenty-seven} \etiq{I'm} \etiq{.}\hfill\trou[66mm]{I'm twenty-seven.}
\item \etiq{that} \etiq{who's} \etiq{?}\hfill\trou[66mm]{Who's that?}
\end{questions}

\exercice[lire,ecrire,niveau=2]{À la place de Sam}{Voici la carte de membre de Sam. Vous êtes Sam :
répondez aux questions de l'employée par une phrase complète.}
\begin{center}
\begin{tikzpicture}
  \fill[otNavy,rounded corners=3mm] (0,0) rectangle (92mm,48mm);
  \fill[white,rounded corners=2mm] (3mm,3mm) rectangle (89mm,35mm);
  \node[anchor=west,text=white,font=\otBold\fontsize{12}{14}\selectfont] at (5mm,41.5mm) {Fit Brooklyn};
  \node[anchor=east,text=otGreen,font=\otSemi\fontsize{7}{8}\selectfont] at (87mm,41.5mm) {\otCaps[10]{Member card}};
  \otDessinAvatar{sam}{(18mm,19mm)}{10mm}{otTealTint}
  \node[anchor=north west,align=left,font=\fontsize{9.4}{12.6}\selectfont,text=otInk] at (32mm,32mm)
    {\enb{\textbf{Name:} Sam Brooks}\\ \enb{\textbf{Age:} 24}\\ \enb{\textbf{Phone:} (718) 555-0162}\\
     \enb{\textbf{Email:} sam.brooks@example.com}\\ \enb{\textbf{Nationality:} American}};
\end{tikzpicture}
\end{center}
\begin{questions}
\item What's your last name? \hfill \trou[62mm]{It's Brooks.}
\item How old are you? \hfill \trou[62mm]{I'm twenty-four.}
\item What's your phone number? \hfill \trou[62mm]{It's 718-555-0162.}
\item Are you British? \hfill \trou[62mm]{No, I'm not. I'm American.}
\end{questions}

\exercice[ecouter,ecrire,niveau=1,piste=3.11]{On épelle}{Écoutez : la voix épelle six noms de
famille. Écrivez-les.}
\begin{questions}[3]
\item \trou[30mm]{Carter}
\item \trou[30mm]{Lopez}
\item \trou[30mm]{Brooks}
\item \trou[30mm]{Hughes}
\item \trou[30mm]{Gray}
\item \trou[30mm]{Quinn}
\end{questions}

\exercice[ecouter,ecrire,niveau=2,piste=3.12]{Numéros et âges}{Quatre personnes donnent leur
numéro de téléphone et leur âge. Complétez le tableau.}
\begin{center}\otTableau
\begin{tabular}{@{}l@{\hspace{14mm}}l@{\hspace{14mm}}l@{}}
\otTH{Prénom} & \otTH{Téléphone} & \otTH{Âge}\\\midrule
{\otSemi 1 · Ana}   & \trou[40mm]{212-555-0134} & \trou[14mm]{31}\\[1mm]
{\otSemi 2 · Marco} & \trou[40mm]{347-555-0119} & \trou[14mm]{45}\\[1mm]
{\otSemi 3 · Yuki}  & \trou[40mm]{646-555-0172} & \trou[14mm]{23}\\[1mm]
{\otSemi 4 · Liam}  & \trou[40mm]{718-555-0156} & \trou[14mm]{38}\\\bottomrule
\end{tabular}\end{center}

\exercice[lire,ecrire,niveau=2]{Des adresses e-mail}{Voici quatre adresses e-mail, telles qu'on
les dit. Écrivez-les.}
\begin{questions}
\item jack dot miller at example dot com \hfill \trou[64mm]{jack.miller@example.com}
\item emma underscore clarke at example dot co dot uk \hfill \trou[64mm]{emma\_clarke@example.co.uk}
\item studio at example dot org \hfill \trou[64mm]{studio@example.org}
\item sam dot brooks twenty-four at example dot com \hfill \trou[64mm]{sam.brooks24@example.com}
\end{questions}

\exercice[ecrire,niveau=2]{Traduisez}{Traduisez en anglais.}
\begin{questions}[1][french]
\item J'ai vingt-sept ans.\lignes{1}
\item Comment ça s'écrit ?\lignes{1}
\item Quel est ton numéro de téléphone ?\lignes{1}
\item Quel âge a Jack ? — Il a trente ans.\lignes{1}
\item C'est une adresse e-mail.\lignes{1}
\item Qui est-ce ? — C'est Emma.\lignes{1}
\item Pardon ? Encore une fois, s'il vous plaît.\lignes{1}
\item Mon nom de famille est Martin.\lignes{1}
\end{questions}

\exercice[ecrire,niveau=3]{Votre fiche}{Répondez en anglais, par une phrase complète, avec vos
propres informations.}
\begin{questions}
\item What's your first name?\lignes{1}
\item What's your last name? How do you spell it?\lignes{1}
\item How old are you?\lignes{1}
\item Where are you from?\lignes{1}
\item What's your phone number?\lignes{1}
\item What's your email address?\lignes{1}
\end{questions}

\exercice[ecouter,parler,niveau=3,piste=3.13]{À vous de parler !}{Vous vous inscrivez à
\en{Fit Brooklyn}. Répondez à l'employée \textbf{à voix haute} pendant chaque pause.}
\begin{center}\otTableau
\begin{tabular}{@{}>{\otSemi\fontsize{8.6}{10}\selectfont\color{otNavy}}l@{\hspace{3mm}}>{\enOn}p{70mm}<{\enOff}@{\hspace{6mm}}p{78mm}@{}}
 & \otTH{L'employée dit\ldots} & \otTH{Vous\ldots}\\\midrule
1. & Hello! What's your name, please? & donnez votre prénom et votre nom.\\
2. & How do you spell your last name? & épelez votre nom de famille.\\
3. & How old are you? & donnez votre âge.\\
4. & What's your phone number? & donnez un numéro (le vôtre, ou 917-555-0148).\\
5. & Thank you. Have a nice day! & remerciez et souhaitez-lui une bonne journée.\\\bottomrule
\end{tabular}\end{center}

%% ---------------------------------------------------------------------
\rubrique{bilan}{Faites le point}

\begin{essentiel}
\begin{minipage}[t]{0.56\linewidth}
\otEssTitre{Les phrases clés}
\renewcommand{\arraystretch}{1.32}
\begin{tabular}{@{}>{\enOn\otSemi\fontsize{9.4}{11}\selectfont}l<{\enOff}@{\hspace{3mm}}>{\fontsize{9.2}{11}\selectfont\color{white!78!otNavy}}l@{}}
How do you spell it?          & Comment ça s'écrit ?\\
What's your phone number?     & Quel est ton numéro ?\\
What's your email address?    & Ton adresse e-mail ?\\
How old are you?              & Quel âge as-tu ?\\
I'm twenty-seven.             & J'ai vingt-sept ans.\\
Who's that? — That's Tom.     & Qui est-ce ? — C'est Tom.\\
Sorry? One more time, please. & Pardon ? Encore une fois.\\
\end{tabular}
\end{minipage}\hfill
\begin{minipage}[t]{0.38\linewidth}
\otEssTitre{Questions}
{\enb{\fontsize{9.4}{12}\selectfont \textbf{What · Where · Who · How · How old} + \textbf{be} + sujet}}
\otEssTitre{L'âge}
{\fontsize{9.2}{11.6}\selectfont avec \en{be} : \en{I'm thirty}, jamais \en{I have…}\par}
\otEssTitre{a / an}
{\fontsize{9.2}{11.6}\selectfont \en{a} + son consonne, \en{an} + son voyelle :\par
\en{a phone, an email, an hour}\par}
\otEssTitre{Prononciation}
{\fontsize{9.2}{11.6}\selectfont \en{E} \phon{iː} · \en{I} \phon{aɪ} · \en{G} \phon{dʒiː} · \en{J} \phon{dʒeɪ}\par}
\end{minipage}
\end{essentiel}

\begin{jesais}
\item épeler mon nom et comprendre un mot épelé ;
\item compter de 0 à 100, et écrire les nombres en lettres ;
\item donner mon âge avec \en{be} : \en{I'm twenty-seven} ;
\item donner mon numéro de téléphone et mon adresse e-mail ;
\item poser des questions avec \en{what}, \en{where}, \en{who}, \en{how} et \en{how old} ;
\item choisir entre \en{a} et \en{an} d'après le son qui suit.
\end{jesais}

\clearpage
\begin{testexpress}
\qcm[b]{« J'ai trente ans » se dit :}{\en{I have thirty years.}}{\en{I'm thirty.}}{\en{I'm thirty years.}}
\qcm[c]{13 se dit :}{\en{thirty}}{\en{three-ten}}{\en{thirteen}}
\qcm[a]{40 s'écrit :}{\en{forty}}{\en{fourty}}{\en{fourteen}}
\qcm[b]{Complétez : \en{\rule[-0.6mm]{9mm}{0.5pt} email address}}{\en{a}}{\en{an}}{\en{the a}}
\qcm[c]{Complétez : \en{\rule[-0.6mm]{9mm}{0.5pt} phone number}}{\en{an}}{\en{the a}}{\en{a}}
\qcm[a]{La lettre \en{E} se prononce :}{\phon{iː}, comme notre « i »}{\phon{eɪ}}{\phon{aɪ}}
\qcm[b]{La lettre \en{I} se prononce :}{\phon{iː}}{\phon{aɪ}, comme « aïe »}{\phon{eɪ}}
\qcm[a]{Dans une adresse e-mail, \en{@} se dit :}{\en{at}}{\en{dot}}{\en{arobase}}
\qcm[a]{« Comment ça s'écrit ? » se dit :}{\en{How do you spell it?}}{\en{How spell you it?}}{\en{How is it spell?}}
\qcm[c]{« Qui est-ce ? » se dit :}{\en{What's that?}}{\en{How's that?}}{\en{Who's that?}}
\end{testexpress}

\noindent{\fontsize{9.6}{12}\selectfont\color{otMuted}Niveau visé (CECRL, niveau A1) : « peut
comprendre des nombres, des prix et l'heure ; peut donner et demander des informations
personnelles simples ». Les corrigés de l'unité sont page~\pageref{cor-u3-1} ; les textes des
pistes d'exercices, page~\pageref{trans-3.11}.\par}
